\subsubsection{Further explicit fourfolds and the exact scope of the method}

Write $D^2_\ell(X)$ for numerical--$\ell$-adic homological equality
with rational coefficients in codimension two. The following transfer
rules give concrete further members of the class, without using the
desired equality as a hypothesis about the target.

\begin{proposition}[Proved; algebraic-correspondence transfer]
\label{B1b-correspondence}
Let $X,Y_1,\ldots,Y_m$ be smooth projective fourfolds over $k$, and
suppose $D^2_\ell(Y_i)$ holds. Suppose algebraic correspondences
$a_i\in\operatorname{CH}^4(X\times Y_i)_\Q$ and
$b_i\in\operatorname{CH}^4(Y_i\times X)_\Q$ satisfy
\[
 \sum_i(b_i)_*(a_i)_*=\operatorname{id}
       \quad\hbox{on }H^4(X_{\bar k},\Q_\ell(2)).
\]
Then $D^2_\ell(X)$ holds. In particular this applies to a generically
finite surjective morphism from any $Y_i$ to $X$.
\end{proposition}
\begin{proof}
If $z$ is numerically trivial on $X$, then $(a_i)_*z$ is numerically
trivial on $Y_i$: its intersection with a complementary cycle $w$
equals the intersection of $z$ with $({}^ta_i)_*w$, by the projection
formula. The hypothesis on $Y_i$ therefore gives
$(a_i)_*\operatorname{cl}_\ell(z)=0$. The displayed identity implies
$\operatorname{cl}_\ell(z)=0$. For $f:Y\to X$ of degree $d>0$, take
$a={}^t\Gamma_f$ and $b=d^{-1}\Gamma_f$ and use $f_*f^*=d$.
This includes inseparable maps: only the nonzero degree in $\Q_\ell$
is used. Homological triviality always implies numerical triviality.
\end{proof}

\begin{lemma}[Proved; smooth blow-ups of fourfolds]
\label{B1b-blowup}
Let $X$ be a smooth projective fourfold over $k$, let $Z\subset X$ be
a smooth closed center of positive codimension on every component,
and let $f:W=\operatorname{Bl}_ZX\to X$.
Then $D^2_\ell(W)$ holds if and only if $D^2_\ell(X)$ holds.
\end{lemma}
\begin{proof}
All cohomology below is geometric; the cycles and operations remain
defined over $k$, so the hypothesis $D^2_\ell(X)$ is used over $k$.
A Cartier center gives an isomorphism, so suppose its codimension
is $c\ge2$, component by component. Let $j:E\hookrightarrow W$
be the exceptional divisor, $q:E=\mathbf P(N_{Z/X})\to Z$, and
$\xi=c_1(\mathcal O_E(1))$. The blow-up decomposition
\cite[Lemma~2.1, p.~103]{ShiodaKatsura1979C} is
\[
 H^4(W,\Q_\ell(2))=
 f^*H^4(X,\Q_\ell(2))\oplus
 \bigoplus_{a=1}^{c-1}
 j_*\bigl(\xi^{a-1}q^*H^{4-2a}(Z,\Q_\ell(2-a))\bigr).
\]
All maps shown are induced by algebraic cycle operations.
Every smooth projective variety of dimension at most two has
numerical--homological equality in every codimension: for divisors
use Matsusaka's numerical--$\tau$ equality, recorded in
\cite[p.~361]{BKleiman1968}, and the vanishing of the rational first
Chern class of an algebraically trivial line bundle; zero-cycles
are measured by degree and codimension-zero classes by components.
This applies to $Z$.

Suppose $D^2_\ell(X)$ and let $z$ be numerically trivial on $W$.
Write its cohomology class in the displayed decomposition as
\[
 \alpha=f^*\alpha_0+
      \sum_{a=1}^{c-1}j_*(\xi^{a-1}q^*\beta_a).
\]
Since $f_*\alpha=\alpha_0$ and $f_*z$ is numerically trivial,
$\alpha_0=0$. The self-intersection identity $j^*j_*=-\xi$ gives,
for $b=c-1,c-2,\ldots,1$,
\[
 -q_*\bigl(\xi^{c-1-b}j^*\alpha\bigr)
 =\beta_b+\sum_{a>b}
       q_*(\xi^{c-1+a-b})\,\beta_a.
\]
The left side is the class of a numerically trivial cycle on $Z$:
proper pushforward, pullback between smooth varieties and multiplication
by $\xi$ preserve numerical triviality by the projection formula.
It is therefore zero. Descending induction gives every $\beta_b=0$;
groups of negative degree are simply zero. Hence $\alpha=0$.
Conversely apply Proposition~\ref{B1b-correspondence} to the degree-one
morphism $f$. No factorization theorem for arbitrary birational maps
in positive characteristic is used here.
\end{proof}

\begin{theorem}[Proved with checked source construction; Fermat cubic fourfolds]
\label{B1b-fermat}
Let $\operatorname{char}k\ne3$ and let
\[
 F=\{z_0^3+\cdots+z_5^3=0\}\subset\mathbf P^5_k.
\]
For every $\ell\ne\operatorname{char}k$, numerical and $\ell$-adic
homological equivalence agree for all cycles on $F$ with rational
coefficients. The same holds for any smooth projective $k$-form of $F$.
Together with Ancona's numerical Hodge theorem for $F$, this places
these fourfolds in the homological-positivity class of
Proposition~\ref{prop:B1b-criterion}.
\end{theorem}
\begin{proof}
We first work over an algebraic closure. Let
$S=\{x_0^3+x_1^3+x_2^3+x_3^3=0\}\subset\mathbf P^3$ and
$E=S\cap\{x_3=0\}$. Both are smooth because $3\ne0$ in $k$.
The cohomology of $S$ is algebraic; here is an explicit verification
valid also in characteristic two. Fix a primitive cube root $\zeta$.
For $a,b\in\Z/3$ put
\[
 L^0_{ab}: x_0+\zeta^a x_1=x_2+\zeta^b x_3=0,
 \qquad
 L^1_{ab}: x_0+\zeta^a x_2=x_1+\zeta^b x_3=0.
\]
These are lines on $S$. Within either family, two distinct lines meet
exactly when their first or their second indices agree; a line
$L^0_{ab}$ meets $L^1_{cd}$ exactly when $a+d=b+c$ in $\Z/3$.
These assertions follow by solving the four linear equations. They
hold in every characteristic different from three: the residual
equation is a difference of two distinct cube roots, which is nonzero.
Distinct meeting lines have intersection number one, and adjunction
gives self-intersection $-1$.

In the order
$L^0_{00},L^0_{01},L^0_{02},L^0_{10},L^0_{11},L^1_{00},L^1_{01}$,
the intersection matrix is
\[
 G=\begin{pmatrix}
 -1&1&1&1&0&1&0\\
 1&-1&1&0&1&0&1\\
 1&1&-1&0&0&0&0\\
 1&0&0&-1&1&0&0\\
 0&1&0&1&-1&1&0\\
 1&0&0&0&1&-1&1\\
 0&1&0&0&0&1&-1
 \end{pmatrix},\qquad \det G=1.
\]
The attached exact program checks this determinant and all 351 pairs
of the 27 coordinate-pairing lines over $\mathbf F_4$.
The preceding cube-root calculation proves the displayed incidence
matrix in every characteristic different from three; the finite program
is not used to infer this characteristic-independent assertion.
Over $\C$, weak Lefschetz gives
$b_1=b_3=0$, and
$c(TS)=(1+H)^4/(1+3H)$ gives
$\int_Sc_2(TS)=\int_S3H^2=9$, hence $b_2=7$.
For each fixed $\ell$, the smooth proper Fermat family over
$\Z[1/(3\ell)]$ and smooth proper base change give the same Betti
numbers in every characteristic considered.
Thus the seven independent line classes span $H^2(S,\Q_\ell(1))$;
$H^0$ and $H^4$ are also algebraic and odd cohomology is zero.
K\"unneth now shows that all the cohomology of $S\times S$ is
algebraic. Poincar\'e duality then gives $D^2_\ell(S\times S)$:
an algebraic class pairing to zero with every algebraic cycle pairs
to zero with the whole complementary cohomology space.

Shioda--Katsura's explicit construction
\cite[Lemmas~1.1--1.2, pp.~98--99; Theorem~1.7, pp.~101--102]
{ShiodaKatsura1979C}, with degree three and $r=s=2$, supplies a morphism
\[
 \psi:\operatorname{Bl}_{E\times E}(S\times S)\longrightarrow F
\]
of generic degree three. Explicitly its rational map before the blow-up is
\[
 ([x],[y])\longmapsto
 [x_0y_3:x_1y_3:x_2y_3:-x_3y_0:-x_3y_1:-x_3y_2].
\]
The sum of the cubes is zero by the two equations of $S$. Its base
ideal is locally $(x_3,y_3)$ along $E\times E$, so blowing up this
smooth codimension-two center resolves the map. On the dense affine
chart $x_3=y_3=1$, an inverse to a general projective target point
requires a scalar $a$ satisfying
$a^3=-1/(z_0^3+z_1^3+z_2^3)$; its three distinct choices give generic
degree three. Properness makes the morphism surjective.
Lemma~\ref{B1b-blowup} proves $D^2_\ell$ for its source, and
Proposition~\ref{B1b-correspondence} proves it for $F$.

Outside middle degree, the even cohomology of a smooth hypersurface
$F$ is generated by powers of its hyperplane class and its odd
cohomology vanishes \cite[Lemma~2.3, p.~103]{ShiodaKatsura1979C}.
This proves equality for all codimensions. Geometric base change and
Lemma~\ref{lem:B-basechange-num} give the statement over $k$, including
forms of $F$.

Finally \cite[Proposition~A.9]{BAncona2021} proves numerical primitive
positivity for this exact Fermat fourfold in every characteristic
different from three. Its algebraic Lefschetz splitting is explicit:
$H^2=\Q_\ell h$, $H^6=\Q_\ell h^3$, $\deg h^4=3$, and
\[
 \alpha=\frac{\deg(\alpha h^2)}{3}h^2+
       \left(\alpha-\frac{\deg(\alpha h^2)}{3}h^2\right),
\]
with the second term primitive. Equality already proved identifies
this numerical primitive form with its homological form. It is
therefore positive definite, as required.
\end{proof}

Consequently an explicit covered class consists of abelian fourfolds,
Fermat cubic fourfolds and their geometric forms, together with all
fourfolds obtained by finitely many applications of
Proposition~\ref{B1b-correspondence} and Lemma~\ref{B1b-blowup}.
The seeds also furnish special fibres for
Proposition~\ref{prop:B1b-scope}, whenever its actual smooth-family
and algebraic-Lefschetz hypotheses hold. An arbitrary fourfold is not
asserted to specialize to one of these seeds with its relevant cycle
data.

The exact domain of the \emph{componentwise definiteness argument}
is intrinsically specified by Proposition~\ref{prop:B1b-criterion}:
the Lefschetz projections must preserve algebraic classes and the
signed primitive intersection form on \emph{homological} algebraic
classes must be positive definite. These are independent cohomological
and positivity hypotheses; numerical--homological equality is the
conclusion, not an assumed property of the target. Under the additional
numerical-positivity hypothesis, the converse in that proposition
identifies the remaining issue precisely with the absence of a
homological radical. The conditional method criterion is complete.
The examples and transfer rules above certify substantial geometric
membership; they are not an assertion that this list exhausts all
geometric fourfolds in that intrinsic class. Ancona's
\cite[Remark~A.10]{BAncona2021} likewise distinguishes his rank-two
method's proved examples from a classification of all fourfolds.
